\section{Balancing all prime-power races in a single multiplicative interval}
\label{sec:balancing}
Fix a positive integer $q$, put $G=(\mathbb Z/q\mathbb Z)^\times$ and
$d=|G|$, and let $\chi$ run over the characters of $G$. Write $\chi^*$
for its primitive character and $m_\chi$ for the order of vanishing of
$L(s,\chi^*)$ at $1/2$. For $a\in G$ define
\begin{align}
 \psi(x;q,a)&=\sum_{n\leq x,\ n\equiv a\pmod q}\Lambda(n),&
 b_a&=\frac{2}{d}\sum_\chi\overline{\chi(a)}m_\chi,\nonumber\\
 D_a(x)&=\frac{\psi(x;q,a)-x/d}{\sqrt x}+b_a,&
 D(x)&=(D_a(x))_{a\in G}.\label{eq:racevec}
\end{align}
The constants $b_a$ are real by conjugation. They incorporate possible
central zeros, so even characters with unknown central nonvanishing
are included. Using the von Mangoldt weight counts prime powers as well as primes.

\paragraph{Fixed-interval balancing theorem.}
Assume the upstream Dirichlet GRH. There exist $C_q>1$ and $r_q>0$
such that, for every sufficiently large real $X$,
\[
 \{y\in\mathbb R^d:\|y\|_2\leq r_q\}
 \ \subseteq\ \operatorname{conv}\{D(n): n\in\mathbb Z,\ X\leq n\leq C_qX\}.
 \tag{B}\label{eq:ballrace}
\]
Consequently every such $y$, including zero, is exactly a convex
combination of at most $d+1$ integer samples from that interval, with
the same sample points and weights for every residue class. Every
unit linear combination of the $d$ coordinates attains values at least
$r_q$ and at most $-r_q$ in the interval. The theorem gives existence
with constants depending on the fixed modulus. It requires neither
simplicity of zeros nor rational independence of their ordinates.

\subsection{A local averaging argument from the Lean proof}

The compact proofs of \texttt{IsPolar.eventually\_moment\_le} and
\texttt{IsPolar.exists\_window\_sub\_mul\_le} use high Laplace moments to
find a one-sided small value on $[t,t^2]$ in log coordinates.
They are in
\href{https://github.com/lsernau/grh-lean/blob/72fead102f663d5c7a407964a13dc0cf9d977207/grh-compact/src/New/Endgame/OneSidedCriterion.lean}{\texttt{New/Endgame/OneSidedCriterion.lean}}.
The following vector version explains how a common choice of positive
weights replaces a separate choice of a point for each direction.

\paragraph{Laplace averaging lemma.}
Let $E:[0,\infty)\to\mathbb R^d$ be measurable, with
$\|E(v)\|\leq M e^{Av}$, $A\geq0$. Suppose its coordinate Laplace
transforms $F(z)=\int_0^\infty E(v)e^{-zv}\,dv$, initially defined
for $\Re z>A$, extend holomorphically to $\Re z>0$ and a neighborhood of
$0$. For every $h>0$ there is $\delta>0$ such that, for all large $T$,
a probability density supported on $[T,(1+h)T]$ has $E$-average of
norm $O(e^{-\delta T})$. The same estimate holds for a convex
combination of at most $d+1$ values in the interval.

\paragraph{Proof.}
Take $c=1+h/2$. Choose $\lambda>A$ sufficiently large as specified
below, set $k=\lfloor\lambda cT\rfloor$, and let $V$ have density
$\lambda^{k+1}v^ke^{-\lambda v}/k!$ on $v>0$. The disk
$|z-\lambda|\leq\lambda$ meets the imaginary axis only at $0$.
Compactness gives $R>\lambda$ with $F$ holomorphic on a neighborhood
of the closed concentric disk of radius $R$. Cauchy's estimate yields
\[
 \|\mathbb E E(V)\|
 =\lambda^{k+1}\|F^{(k)}(\lambda)\|/k!
 \ll\lambda(\lambda/R)^k.
\]
For the tails, exponential tilting gives
\[
 \mathbb E[e^{AV};\ V\notin I]
 =\left(\frac{\lambda}{\lambda-A}\right)^{k+1}
      \Pr(W\notin I),\qquad I=[T,(1+h)T],
\]
where $W$ has the same shape $k+1$ and rate $\lambda-A$. Its mean
divided by $T$ tends to $c\lambda/(\lambda-A)$.
For all sufficiently large $\lambda$ this limit stays a fixed
positive distance from both endpoints $1,1+h$. The Chernoff bound for
a gamma variable at $u$ times its mean is
$\exp[-(k+1)(u-1-\log u)]$ for the corresponding tail. This follows
by optimizing its moment generating function.
The two endpoint ratios stay bounded away from $1$, so the probability
is at most $2e^{-\kappa_h(k+1)}$ with $\kappa_h>0$ independent of all
sufficiently large $\lambda,T$. Choose $\lambda$ large enough that
$\log(\lambda/(\lambda-A))<\kappa_h/2$. Truncating the density to $I$
thus changes its vector integral by $O(e^{-\delta T})$, and its mass
is $1-O(e^{-\delta T})$, after decreasing $\delta$ if necessary.
The normalized integral belongs to the closed convex hull. Approximate
it to $e^{-\delta T}$ by a finite convex combination and apply
Carath\'eodory's theorem.\hfill$\square$

For the character coordinate of \eqref{eq:racevec}, write
$\psi(x,\chi)=\sum_{n\le x}\chi(n)\Lambda(n)$,
$e_\chi=1$ for the principal character and $0$ otherwise, and
$s=z+1/2$. The Laplace transform of
$e^{-v/2}(\psi(e^v,\chi)-e_\chi e^v)+2m_\chi$ is
\[
 -\frac{1}{s}\frac{L'}{L}(s,\chi)-\frac{e_\chi}{s-1}
                +\frac{2m_\chi}{z}.
\]
The pole at $s=1$ cancels for the principal character. The residue
at $z=0$ before adding the final term is $-2m_\chi$. Finite Euler
factors removed from primitive $L$-functions have zeros only on
$\Re s=0$. Dirichlet GRH therefore gives the continuation required
by the lemma. Fourier inversion gives all the coordinates in
\eqref{eq:racevec}. The lemma already gives exponentially accurate
balance in each log-window $[T,(1+h)T]$.
For exact balance in a log-window of fixed length we use the
individual nonzero frequencies.

\subsection{Simultaneous distinction of the zeros}

\paragraph{Multiplicity rank theorem.}
For each fixed $q$ there is $K_q$ such that if $T$ is sufficiently
large and $K_q\log T\leq H\leq T$, the vectors
\[
 M(\rho)=\big(\operatorname{ord}_\rho L(s,\chi^*)\big)_{\chi\bmod q}
 \quad (0<\Re\rho<1,\ T<\Im\rho<T+H)
\]
span $\mathbb R^d$. Distinct zeros index the vectors; a zero absent
from one $L$-function contributes a zero entry. This assertion about
complex zeros holds unconditionally. Under Dirichlet GRH the same
vectors are indexed by distinct positive ordinates.

\paragraph{Proof.}
Choose a prime $p_a>2$ in each class $a\in G$ by Dirichlet's theorem.
Banks's short-interval Landau formula gives
\[
 A_{a\chi}:=\sum_{T<\Im\rho<T+H}p_a^\rho M(\rho)_\chi
  =-\frac{H\log p_a}{2\pi}\chi(a)+O_q(\log T).
\]
Here the sum is over the union of distinct nontrivial zeros. At these
primes $\chi^*(p_a)=\chi(a)$. The implied constant can be fixed for
the finite list of characters and chosen primes. The row-scaled
character table multiplying $H$ is invertible. Choose $K_q$ so its
smallest singular value times $H$ exceeds the operator norm of the
error matrix. Then $A$ is invertible. Its rows belong to the complex
span of the real vectors $M(\rho)$, proving their full real rank.
The input is Theorem 1.2 of W.~Banks,
\href{https://arxiv.org/abs/2302.07073}{\emph{Comparing zeros of distinct
Dirichlet $L$-functions}}. The conductor-one case is the classical
Landau formula in the same paper.\hfill$\square$

This yields a nonsingular integer $d$-by-$d$ multiplicity matrix in
every such height window, including when different characters share zeros.
In particular one can choose $\gg_q T/\log T$ pairwise disjoint
bases with ordinates in $[T,2T]$, by dividing it into intervals of
length $2K_q\log(2T)$. Under Dirichlet GRH these are bases of
multiplicity vectors attached to distinct heights. 
Banks also obtains bounds uniform in the conductor for a pair of
characters. The rank argument here keeps $q$ fixed. The entries of the
matrix record multiplicities. Its nonsingularity places no condition
on rational linear relations among the heights themselves.

\subsection{Isolating oscillations with positive weights}

For a nonzero ordinate $\gamma$ in the union of the spectra put
$\rho_\gamma=1/2+i\gamma$ and
\[
 c_{\gamma,a}=-\frac{1}{d\rho_\gamma}
               \sum_\chi \overline{\chi(a)} M(\rho_\gamma)_\chi,
 \qquad c_\gamma=(c_{\gamma,a})_{a\in G}.
\]
The finite Fourier transform and the rank theorem show that the
vectors $c_\gamma$ for positive $\gamma$ span $\mathbb C^d$.
Thus the real and imaginary parts of finitely many of them span $\mathbb R^d$.

On $v>1$, the explicit formula, as an identity of distributions, reads
\[
 D(e^v)=\sum_{\gamma\ne0}c_\gamma e^{i\gamma v}
          +O_q((1+v)e^{-v/2}).\tag{E}\label{eq:explicitvec}
\]
To see the error term, first use a primitive character. The main pole
gives $e_\chi e^v$, each nontrivial zero contributes $-e^{\rho v}/\rho$
with multiplicity, and the central contribution is canceled by
\eqref{eq:racevec}. The pole of the Perron kernel at zero and a possible
trivial zero there contribute $O_q(1+v)$ before normalization.
The other trivial-zero terms form a geometrically decreasing series
for $v>1$. Replacing a primitive character by the one modulo $q$
changes $\psi$ by $O_q(v+1)$, since only powers of primes dividing $q$
are affected. The choice at a jump is immaterial to integration.
A primary reference for the boundary-value formula and its convergence
for smooth averages is Kaczorowski and Ramar\'e, Theorem 1(d),(e),
\href{https://doi.org/10.4064/aa106-3-6}{\emph{Almost periodicity of some
error terms in prime number theory}}.

Take a smooth probability density $w$, positive on $(0,1)$ and supported
on $[0,1]$, and let $w_n$ be its $n$-fold convolution, supported on
$[0,n]$. Put $\widehat w(t)=\int w(u)e^{itu}\,du$. For each nonzero
real $t$, $|\widehat w(t)|<1$, and $\widehat w$ decreases rapidly.
For a fixed $\omega$, integration of (E) shows
\begin{align*}
 &\int_0^n w_n(u) D(e^{V+u})e^{-i\omega(V+u)}\,du \\
 &\qquad=\sum_{\gamma\ne0}
      c_\gamma e^{i(\gamma-\omega)V}\widehat w(\gamma-\omega)^n
       + O_q((1+V)e^{-V/2})\qquad(V\ge 2).
\end{align*}
The error is uniform in $n$ since $(1+v)e^{-v/2}$ decreases for $v\ge2$.
The sum converges absolutely for $n\ge1$ by rapid decay and the
classical $O_q(U\log(2U))$ zero count up to height $U$. The union of
the spectra is locally finite. Hence for each fixed $\omega$ the
supremum of $|\widehat w(\gamma-\omega)|$ over $\gamma\ne\omega$ is
strictly below $1$. All off-frequency terms tend uniformly in $V$ to
zero as $n\to\infty$. At $\omega=0$ the sum therefore tends to zero.
At a selected $\omega=\gamma$ it tends to $c_\gamma$.
For the arithmetic remainder take first a sufficiently large fixed
$n$, then $V$ sufficiently large.

Average $D(e^v)$ against the nonnegative densities proportional to
$w_n(v-V)(1\pm\cos(\gamma v))$ and
$w_n(v-V)(1\pm\sin(\gamma v))$. Their masses tend uniformly to $1$
because $|\widehat w(\gamma)|<1$. The averages approach, to any
prescribed accuracy, the vertices $\pm\Re c_\gamma$ and
$\pm\Im c_\gamma$ for the selected finite set of positive ordinates.
Here the opposite signs for sine merely interchange the two vertices.
The convex hull of these ideal vertices contains a ball about the
origin because their real span has full dimension. Concretely, if its
inradius is $r>0$, approximating every vertex within $r/4$ preserves
a ball of radius $3r/4$, by the support-function characterization of
convex inclusion.

Finally replace $D(e^v)$ in the averages by $D(\lfloor e^v\rfloor)$.
Between consecutive integers the numerator's arithmetic sum is
constant, and $\psi(x;q,a)\leq x\log(2x)$. The mean value theorem
bounds the change by $O_q((1+v)e^{-v/2})$, uniformly for large $v$.
The averages now belong to the convex hull of a \emph{finite} set of
integer samples. Set $V=\log(2X)$, so those integers lie in
$[X,2e^nX]$. Taking $C_q=2e^n$ and $r_q=r/2$ proves
\eqref{eq:ballrace}. Carath\'eodory's theorem gives the sample bound.
\hfill$\square$

\paragraph{Relation to the open prime race problem.}
A ball in the convex hull allows thresholds to be mixed. A prescribed
ranking of raw prime counts asks for all of its inequalities to hold
at one threshold. Even a circle has a convex hull containing points
that the circle never visits. For results on the single-threshold
problem under hypotheses on the zero ordinates, see G.~Martin and
N.~Ng, \href{https://doi.org/10.1090/tran/7996}{\emph{Inclusive prime
number races}}. The conditional almost-periodicity theory of Kaczorowski and Ramar\'e
(Theorem 1(e)) gives a second perspective. Their corrected vector has
Stepanov almost-periodic coordinates, zero mean after removing the
central terms, and the nonzero Fourier coefficients displayed in (E),
up to an overall sign. The rank theorem makes every nonzero real
projection nonconstant. Positive and negative values of a nonzero
mean-zero projection then recur on sets of positive measure in
intervals of bounded log-length: an interval on which it exceeds a
positive level on a positive-measure subset can be repeated using
relatively dense Stepanov almost-periods; the same applies to its
negative. The positive compact averages above give a direct proof
uniform across directions. This classical theory and Banks's
explicit formula account for the analytic inputs. We have not
established priority for the rank or convex-hull formulations.

\paragraph{An exact finite example.}
Take $q=4$. Its two central orders vanish. The nonprincipal central
value is a positive alternating series. For the principal character,
$\zeta(1/2)=\eta(1/2)/(1-\sqrt2)$ and the alternating eta series is
positive. The three integers $293,283,228$ lie in $[200,400]$.
For $v(n)=(\psi(n;4,1)-n/2,\psi(n;4,3)-n/2)$ the certified lower
bounds for the cyclic determinants are
\[
 \det(v(293),v(283))\ge5.625,\quad
 \det(v(283),v(228))\ge41.693,\quad
 \det(v(228),v(293))\ge14.050.
\]
All three are positive. The identity
$\det(v_2,v_3)v_1+\det(v_3,v_1)v_2+\det(v_1,v_2)v_3=0$
therefore gives strictly positive balancing weights. On replacing
$v(n)$ by $D(n)=v(n)/\sqrt n$, multiply the associated weight by
$\sqrt n$ and renormalize. The file
\texttt{checks/prime\_race\_triangle.py} verifies the three bounds using
integer sieving and rational upper and lower estimates for logarithms.
This example is independent of GRH.
